\section{Orquestación con Kubernetes}
\label{sec:kubernetes}

Kubernetes~\cite{kubernetes_burns2016} es un sistema de orquestación de contenedores que gestiona de forma
declarativa el ciclo de vida de los servicios: arranque, reinicio ante fallos y escalado.
A diferencia de Docker Compose, que está pensado para un único equipo, Kubernetes puede
ejecutarse sobre infraestructura en la nube sin modificar el código de la aplicación,
lo que lo convierte en el entorno de producción natural para proyectos de investigación
como CyberNEMO. Para las pruebas locales se utilizó Minikube~\cite{minikube_docs}, que levanta un
clúster Kubernetes completo en el propio equipo de desarrollo sin necesidad de
infraestructura en la nube.

La unidad mínima de despliegue en Kubernetes es el \textit{pod}. Cuando varios contenedores
se agrupan en el mismo pod, comparten automáticamente la interfaz de red, es decir, todos
ven el mismo \texttt{localhost}. Este comportamiento resuelve de forma nativa el mismo
problema de aislamiento de red que en Docker se solucionó con el espacio de red compartido.
El pod principal, llamado \texttt{studio}, agrupa voctocore, todas las fuentes de vídeo
(cam1--4, break, intro, \textit{stream blanker} y \textit{audio manager}) y el servicio de telemetría. Al compartir pod, todas las conexiones
TCP internas funcionan sin modificar una sola línea de código. RabbitMQ se despliega en un
pod independiente; la telemetría le publica sus eventos mediante \acs{AMQP}.

El despliegue se define en cinco ficheros \ac{YAML}:
\texttt{studio.yaml} define el pod principal con todos sus contenedores y los puertos
expuestos; \texttt{rabbitmq.yaml} despliega RabbitMQ como pod independiente;
\texttt{configmap.yaml} centraliza todos los parámetros de configuración (resolución,
puertos, fuentes de vídeo) como variables de entorno accesibles por todos los contenedores,
de forma que un solo cambio aquí se aplica a todo el clúster; \texttt{pvc.yaml} reserva
1~GiB de almacenamiento persistente para los registros de telemetría, que así sobreviven
al reinicio del pod; y \texttt{secret.yaml} almacena las credenciales de RabbitMQ de
forma segura, sin escribirlas directamente en los manifiestos. A diferencia del escenario
Docker, donde la imagen se construye en el propio equipo, Kubernetes descarga la imagen
directamente de Docker Hub (\texttt{martin1210235/voctomix:latest}), lo que permite
desplegar el sistema en cualquier nodo del clúster sin necesidad de construirla localmente.
Los manifiestos completos se reproducen en la sección~\ref{sec:anexoC6} del Anexo~C.

El orden de arranque también requiere coordinación. Si el pod \texttt{studio} intentase
arrancar antes de que RabbitMQ estuviera listo, la telemetría fallaría al intentar
conectarse. Para evitarlo, un pequeño contenedor auxiliar se ejecuta primero y simplemente
espera hasta que RabbitMQ responde; solo entonces arranca el resto del pod. Del mismo modo,
una comprobación periódica sobre el puerto~9999 confirma que voctocore está operativo antes
de que el sistema empiece a aceptar conexiones externas.

Para arrancar el sistema, el operador lanza en PC1 un comando que primero comprueba
si Minikube ya está en marcha y, si no, lo inicia. A continuación aplica todos los
manifiestos YAML al clúster y espera a que los pods estén completamente listos. Una vez
que el sistema está operativo, establece un reenvío de puertos
(\texttt{kubectl port-forward --address 0.0.0.0}) que hace accesibles desde la red local
los puertos de voctocore (control, \textit{previews}, salida de programa). En PC2, el operador
lanza voctogui apuntando a la IP de PC1, exactamente igual que en el escenario de dos PCs
nativo. Para detener el sistema basta con cerrar el terminal de PC1, lo que cierra
automáticamente el reenvío de puertos; el clúster y los pods permanecen activos y pueden
retomarse en cualquier momento sin necesidad de volver a desplegarlos (figura~\ref{fig:escenario_kubernetes}).

\begin{figure}[H]
  \centering
  \includegraphics[width=\textwidth]{figuras/escenario_kubernetes.png}
  \caption{Escenario de despliegue con Kubernetes.}
  \label{fig:escenario_kubernetes}
\end{figure}

Los cuatro escenarios de despliegue cubren un rango amplio de necesidades: el escenario
de un PC es el punto de entrada natural para el desarrollo y las pruebas; el de dos PCs
permite separar el servidor del puesto del realizador en una red local; Docker elimina
la instalación de dependencias y hace el sistema portable a cualquier máquina; y
Kubernetes lo integra en infraestructuras de producción como la del proyecto CyberNEMO.
En todos los casos el código de Voctomix~2.0 es idéntico, y solo cambia el entorno
de ejecución.
